\section{Mobile-First dan Breakpoint}

Pendekatan \textbf{mobile-first} menulis CSS dasar untuk layar kecil terlebih dahulu, kemudian menambahkan aturan untuk layar lebih besar menggunakan \code{@media (min-width: ...)}. Ini adalah kebalikan dari pendekatan desktop-first yang menggunakan \code{max-width}. Mobile-first menghasilkan CSS yang lebih ringan karena mobile hanya memuat aturan dasar, sementara aturan tambahan hanya ditambahkan untuk layar yang lebih besar \cite{webdevCss}.

Breakpoint adalah titik lebar viewport di mana layout berubah. Breakpoint umum: 480px (mobile), 768px (tablet), 1024px (desktop kecil), 1200px (desktop besar). Namun, breakpoint sebaiknya dipilih berdasarkan kapan konten membutuhkan penyesuaian, bukan berdasarkan perangkat tertentu. Dengan Flexbox/Grid yang responsif secara alamiah, jumlah breakpoint yang diperlukan sering lebih sedikit dari yang diperkirakan \cite{mdnCss}.

\begin{webcode}[language=CSS, caption={Pendekatan mobile-first}]
/* Mobile (base) */
.card-grid {
  display: grid;
  grid-template-columns: 1fr;
  gap: 1rem;
  padding: 1rem;
}

/* Tablet ke atas */
@media (min-width: 768px) {
  .card-grid { grid-template-columns: repeat(2, 1fr); }
}

/* Desktop ke atas */
@media (min-width: 1024px) {
  .card-grid { grid-template-columns: repeat(3, 1fr); }
  .sidebar { display: block; width: 250px; }
}
\end{webcode}

\begin{konsep}
\textbf{Mobile-First vs Desktop-First:}
\begin{itemize}
  \item \textbf{Mobile-first} (\code{min-width}): CSS base untuk mobile, lalu tambah fitur di layar besar. Lebih ringan di mobile; direkomendasikan.
  \item \textbf{Desktop-first} (\code{max-width}): CSS base untuk desktop, lalu kurangi fitur di layar kecil. Lebih intuitif saat mendesain dari mockup desktop, tapi menghasilkan CSS lebih berat di mobile.
  \item Kedua pendekatan bisa dicampur, tetapi konsistensi meningkatkan maintainability \cite{cssTricks}.
\end{itemize}
\end{konsep}

